\documentclass[12pt,a4paper]{article}
\usepackage[a4paper,margin=2.5cm,headheight=15pt]{geometry}
\usepackage[T1]{fontenc}
\usepackage[utf8]{inputenc}
\usepackage{microtype}
\usepackage{setspace}
\usepackage{parskip}
\usepackage{fancyhdr}
\usepackage{titlesec}
\usepackage{array}
\usepackage{booktabs}
\usepackage{tabularx}
\usepackage{enumitem}
\usepackage{xcolor}
\usepackage{listings}
\usepackage[hidelinks]{hyperref}
\usepackage{xurl}
\usepackage[backend=biber,style=numeric,sorting=none]{biblatex}
\addbibresource{references.bib}

\setstretch{1.15}
\setlength{\parindent}{0pt}
\setlength{\parskip}{0.55em}
\setlist[itemize]{leftmargin=1.5em,itemsep=0.2em,topsep=0.3em}
\setlist[enumerate]{leftmargin=1.8em,itemsep=0.25em,topsep=0.3em}
\renewcommand{\arraystretch}{1.2}
\pagestyle{fancy}
\fancyhf{}
\lhead{DRY and SRP}
\rhead{Software Design Principles}
\cfoot{\thepage}
\titleformat{\section}{\Large\bfseries}{\thesection.}{0.6em}{}
\titleformat{\subsection}{\large\bfseries}{\thesubsection}{0.6em}{}
\titleformat{\subsubsection}{\normalsize\bfseries}{\thesubsubsection}{0.6em}{}

\definecolor{codeblue}{RGB}{28,75,135}
\definecolor{codegray}{RGB}{95,95,95}
\definecolor{codegreen}{RGB}{35,115,70}
\definecolor{codebg}{RGB}{248,249,251}
\lstdefinestyle{cppstyle}{
  language=C++,
  basicstyle=\ttfamily\small,
  keywordstyle=\color{codeblue}\bfseries,
  commentstyle=\color{codegreen}\itshape,
  stringstyle=\color{red!65!black},
  numbers=left,
  numberstyle=\tiny\color{codegray},
  stepnumber=1,
  numbersep=8pt,
  backgroundcolor=\color{codebg},
  frame=single,
  rulecolor=\color{black!18},
  breaklines=true,
  breakatwhitespace=false,
  showstringspaces=false,
  tabsize=4,
  columns=fullflexible,
  keepspaces=true,
  captionpos=b
}
\lstset{style=cppstyle}

\begin{document}
\begin{titlepage}
  \centering
  \vspace*{2cm}
  {\Large SOFTWARE DESIGN PRINCIPLES\par}
  \vspace{1.3cm}
  {\huge\bfseries DRY and SRP:\par}
  \vspace{0.35cm}
  {\LARGE\bfseries Applying the Principles Appropriately\par}
  \vspace{1cm}
  {\large Research Report\par}
  \vfill
  \begin{tabular}{@{}ll@{}}
    \textbf{Student:} & \rule{9cm}{0.4pt}\\[0.8cm]
    \textbf{Course:} & \rule{9cm}{0.4pt}\\[0.8cm]
    \textbf{Language of examples:} & C++17\\[0.8cm]
    \textbf{Date:} & 13 November 2026
  \end{tabular}
  \vfill
  {\large\itshape A study of duplication, responsibility, cohesion, coupling, and design trade-offs\par}
\end{titlepage}
\pagenumbering{roman}
\tableofcontents
\newpage
\pagenumbering{arabic}
\begin{quote}
\textbf{Scope note:} This report follows the supplied assignment brief. The code samples in this repository are small, self-contained teaching examples designed to make the design decisions visible. They are not claimed to be excerpts from a production system. The report distinguishes established definitions from the analysis and decisions made for these examples.
\end{quote}

\begin{abstract}

DRY (Don't Repeat Yourself) and SRP (Single Responsibility Principle) are widely used principles for organizing maintainable software. DRY is concerned with avoiding duplicated knowledge, while SRP is concerned with grouping behavior according to the reasons that cause it to change. Neither principle means that all similar-looking code must be shared or that every method must belong to a separate class. Applied without context, DRY can produce premature abstractions and unwanted dependencies; SRP can produce excessive classes, interfaces, wrappers, and indirection.

This report explains the purpose and practical application of both principles, examines situations in which their strict or superficial application can make a design worse, and evaluates five code-review scenarios. The examples use C++17 to compare duplicated and shared validation, independent business rules, an over-generalized percentage abstraction, a multi-purpose order service, a cohesive student record, and excessive getter/setter classes. The central conclusion is that design quality should be judged by the stability of the concepts being represented, the reasons components change, cohesion, coupling, and the cost of future maintenance—not by the number of duplicated lines or classes.

\end{abstract}

\section{Introduction}

Software rarely remains unchanged after its first implementation. Business policies evolve, users request new behavior, external services change, and defects reveal assumptions that were not understood at the beginning. A design that is easy to write today can become expensive to modify later when knowledge is scattered across the codebase or when unrelated responsibilities are concentrated in one component.

Two principles often used to address these problems are DRY and SRP:

\begin{itemize}

\item \textbf{DRY} asks developers to avoid maintaining multiple, potentially inconsistent representations of the same knowledge.

\item \textbf{SRP} asks developers to organize behavior so that a component has one coherent responsibility, or one principal reason to change.

\end{itemize}

Both principles aim to make change safer and less expensive. However, neither should be treated as a mechanical rule. Code that looks similar may represent different business policies. A class with several methods may still have one coherent responsibility. Conversely, a design with very little repeated code and many small classes may be harder to understand than a simpler design.

\subsection{Objectives}

This report has four objectives:

\begin{enumerate}

\item Explain the basic meaning and purpose of DRY and SRP.

\item Demonstrate how the principles can improve design problems through representative C++ examples.

\item Analyze over-application, premature abstraction, cohesion, coupling, and the trade-offs involved in refactoring.

\item Make reasoned decisions about five code-review situations instead of automatically refactoring every apparent violation.

\end{enumerate}

\subsection{Scope and approach}

The report uses the assignment brief as its structure and draws on established software-design references. Hunt and Thomas provide the foundational explanation of DRY as a principle about knowledge duplication \cite{thomas2019pragmatic}. Robert C. Martin explains SRP in terms of grouping things that change for the same reason and separating those that change for different reasons \cite{martin2010srp}. Martin Fowler provides a definition of refactoring as behavior-preserving structural change \cite{fowler2004refactoring}, and discusses why coupling must be controlled rather than eliminated indiscriminately \cite{fowler2001coupling}.

The C++ files in the source directory are deliberately small enough to compile independently. They isolate a design idea rather than implement production-grade validation, persistence, email delivery, or student administration. For example, the customer email check is intentionally minimal; a real application would require a clearly specified validation policy and appropriate tests.

\section{Fundamental Concepts}

\section{DRY: Don't Repeat Yourself}

\subsection{Meaning}

DRY stands for \textbf{Don't Repeat Yourself}. In The Pragmatic Programmer, Andrew Hunt and David Thomas formulate the principle around the idea that each piece of knowledge should have one unambiguous, authoritative representation within a system \cite{thomas2019pragmatic}.

This definition is more precise than the common interpretation that programmers should never write the same lines twice. DRY is not primarily a rule about visual similarity. It is a rule about the duplication of \textbf{knowledge, intent, or policy}.

Consider a system in which two operations validate a customer according to the same business policy. If each operation independently implements the same rules, the system stores one policy in two places. A future change to that policy must be applied consistently to both copies. If one copy is missed, the operations no longer agree about what makes a customer valid.

By contrast, two functions may contain similar arithmetic while implementing different business rules. If those rules have different owners or change for different reasons, making them share one abstraction may create an unwanted dependency. The code may look duplicated without representing duplicated knowledge.

\subsection{Problems DRY is intended to prevent}

Duplicated knowledge creates several maintenance risks:

\begin{enumerate}

\item \textbf{Inconsistent updates.} A developer changes one copy but forgets another.

\item \textbf{Behavioral drift.} Initially identical implementations gradually acquire different edge cases or error handling.

\item \textbf{Higher change cost.} A developer must locate every representation of a rule before safely modifying it.

\item \textbf{More testing work.} Multiple implementations of one policy must be checked for consistency.

\item \textbf{Unclear authority.} It becomes difficult to determine which copy represents the current requirement.

\end{enumerate}

These risks matter because requirements are not stable. The original understanding of a rule may change, and the code must continue to express that updated understanding. Hunt and Thomas emphasize that maintenance is not restricted to the period after release; requirements and understanding can change throughout development \cite{thomas2019pragmatic}.

\subsection{Forms of repetition in software}

Duplication can appear in more than one form:

\begin{table}[htbp]
\centering
\small
\begin{tabularx}{\textwidth}{>{\raggedright\arraybackslash}p{0.22\textwidth}XX}
\toprule
Form of duplication & Example & Possible consequence \\
\midrule
Source-code duplication & The same validation condition appears in two functions & One implementation may be updated while the other is missed \\
Business-rule duplication & The same eligibility rule is independently encoded in multiple modules & Modules can disagree about a single policy \\
Data or configuration duplication & A timeout, tax rate, or status mapping is copied into several files & Values can become inconsistent \\
Process duplication & Several services independently implement the same required sequence & A process change must be repeated and coordinated \\
Documentation duplication & The same configuration rule is explained in multiple documents & Documentation becomes contradictory \\
Schema or interface duplication & The same data structure is manually defined in multiple systems & Changes to the structure must be synchronized \\
\bottomrule
\end{tabularx}
\end{table}

Not every form can or should be removed by creating a shared function. For example, client and server code may both need a representation of the same data contract because they run in different environments. In that situation, a shared schema or code-generation process may be more appropriate than forcing both sides to call the same function.

\subsection{Why repeated logic becomes dangerous when requirements change}

Suppose a customer-registration function and a customer-update function both check that:

\begin{itemize}

\item the customer's name is not empty;

\item the email field is not empty; and

\item the email contains an at-sign.

\end{itemize}

At first, the duplicated logic produces the same result. Later, the project changes its policy—for example, it introduces a centralized email-validation library or changes which customer fields are mandatory.

If the developer updates only the registration function, registration and update no longer enforce the same policy. This inconsistency may lead to confusing user experiences and invalid data entering the system through one operation.

The core problem is not that the condition occupies several lines. The problem is that a single policy has multiple independently maintained representations.

\subsection{Common techniques for reducing duplication}

The appropriate technique depends on what is duplicated:

\begin{itemize}

\item \textbf{Extract a function.} Use this when the same small algorithm or condition is needed in multiple places and has a meaningful name.

\item \textbf{Create a policy or domain component.} Use this when several operations must follow the same business rule.

\item \textbf{Centralize configuration.} Store a shared setting in one authoritative configuration source rather than copying its value into several modules.

\item \textbf{Use a schema or code generator.} Use this when multiple representations must be produced from one data definition.

\item \textbf{Share tests or test data.} Common test cases can help verify that multiple entry points implement the same contract.

\item \textbf{Generate documentation where appropriate.} When the same information must appear in several formats, generating those formats from a maintained source can reduce manual synchronization.

\end{itemize}

The goal is not simply to reduce line count. The extracted abstraction should communicate the shared concept, remain easy to understand, and avoid connecting components that should evolve independently.

\subsection{Example: duplicated customer validation}

The repository contains two implementations:

\begin{itemize}

\item \path{src/dry/validation_before.cpp} duplicates the same validation condition in registerCustomer() and updateCustomer().

\item \path{src/dry/validation_after.cpp} extracts the shared policy into isValidCustomer().

\end{itemize}

The second version is preferable in this scenario because both operations intentionally enforce the same customer policy. If the policy changes, the shared function can be updated once. It can also be tested independently.

The example does not claim that checking for an at-sign is sufficient email validation. It uses a deliberately simple condition so that the design issue is visible. A real project should define validation requirements carefully and test them against the cases the product actually supports.

\textbf{Design decision:} share the validation rule because the two operations represent the same policy, not merely because their code happens to look alike.

\section{SRP: Single Responsibility Principle}

\subsection{Meaning}

SRP stands for \textbf{Single Responsibility Principle}. Robert C. Martin describes its central idea as: gather together things that change for the same reason and separate things that change for different reasons \cite{martin2010srp}.

The phrase \textbf{one reason to change} is a practical way to express this idea. A class or module should not be forced to change because of several unrelated sources of requirements.

For example, a component that contains business calculations, report formatting, and database persistence may be affected by three independent groups of requests. A change in the database schema should not normally require modifying the business-calculation logic. Likewise, changing the report layout should not require changing how business rules are calculated.

SRP therefore concerns the organization of responsibilities and change boundaries, not the number of methods in a class.

\subsection{What is a responsibility?}

A responsibility is a coherent area of behavior or knowledge for which a component is accountable. It is often associated with a stakeholder, user group, policy, or technical concern that can request a change.

Examples include:

\begin{itemize}

\item calculating an order total according to pricing rules;

\item storing and retrieving orders;

\item sending an order-confirmation message;

\item validating an order against business constraints; and

\item formatting a student report.

\end{itemize}

A responsibility is not simply one verb or one method. For instance, a student-record component may need methods to update a student's name, add a grade, calculate an average, and determine academic standing. Those methods can form one coherent responsibility if they all maintain or expose the state and rules of a student record.

\subsection{Meaning of "one reason to change"}

The phrase does not mean a class can never be modified for more than one small code-level reason. It means that the class should represent one coherent responsibility rather than several independent concerns.

Consider an OrderService that:

\begin{enumerate}

\item validates an order;

\item calculates the order total;

\item writes the order to storage; and

\item sends a confirmation email.

\end{enumerate}

These actions can change for different reasons:

\begin{table}[htbp]
\centering
\small
\begin{tabularx}{\textwidth}{>{\raggedright\arraybackslash}p{0.27\textwidth}X}
\toprule
Concern & Example change request \\
\midrule
Validation & The business changes which orders are accepted \\
Calculation & Shipping or discount rules change \\
Persistence & The database schema or storage provider changes \\
Notification & The email provider or message template changes \\
\bottomrule
\end{tabularx}
\end{table}

When all of these concerns are implemented in one class, the class becomes a point where unrelated changes accumulate. Its callers may also become exposed to implementation changes that they do not need to know about.

\subsection{How to recognize a class with unrelated responsibilities}

The following questions can reveal a possible SRP problem:

\begin{itemize}

\item Can different stakeholders request changes to different groups of methods?

\item Do some methods implement business policy while others handle storage, networking, or presentation?

\item Does changing one concern require retesting or modifying unrelated behavior?

\item Do the methods form distinct groups with little interaction?

\item Does the class depend on several unrelated technical systems?

\item Does its name hide multiple distinct responsibilities?

\end{itemize}

These are diagnostic signals, not automatic proof. A class with many methods is not necessarily poorly designed, and a class with few methods is not necessarily well designed. The actual responsibilities and likely change patterns matter more than the method count.

\subsection{How responsibilities can be separated}

A class with multiple responsibilities can often be refactored into components with explicit roles:

\begin{itemize}

\item OrderValidator checks whether an order satisfies the rules.

\item OrderCalculator calculates the total.

\item OrderRepository stores the order.

\item ConfirmationSender sends the notification.

\item OrderService coordinates the workflow.

\end{itemize}

The service may still call several methods, but its responsibility is now orchestration. It does not need to own the internal details of validation, calculation, persistence, or email delivery.

This separation can make each component easier to test or replace. However, it is useful only when the new boundaries represent meaningful responsibilities. Splitting code into classes that merely forward calls adds indirection without necessarily improving the design.

\subsection{Example: refactoring an order service}

The repository compares:

\begin{itemize}

\item \path{src/srp/order_service_before.cpp}, where validation, calculation, persistence, and notification are placed in one class.

\item \path{src/srp/order_service_after.cpp}, where those responsibilities are assigned to separate collaborators and the service coordinates them.

\end{itemize}

The refactored version makes it easier to change one concern independently. For example, a real repository implementation could replace the current output-only repository without rewriting the validation rules. Similarly, a different notification mechanism could replace the sender.

These examples simulate persistence and notification with console output. They demonstrate responsibility boundaries, not a functioning database or email integration.

\textbf{Design decision:} split the independent concerns, but keep an orchestration service because coordinating the order-placement process is itself a coherent responsibility.

\section{Advanced Analysis: Misuse, Over-Application, and Trade-Offs}

\section{DRY: Is less duplicated code always better?}

No. A smaller amount of repeated text does not automatically mean a better design. The important question is whether the repeated code expresses the same underlying knowledge.

Two code fragments may look alike for different reasons. A percentage calculation, for example, may be used by both a tax module and a customer-discount module. The arithmetic can be identical while the business policies are independent.

\subsection{Shared code can create unwanted dependencies}

A shared component is a dependency for every component that uses it. This is often desirable when the shared component represents a stable common rule. It can be harmful when it combines concerns that should evolve separately.

Potential problems include:

\begin{itemize}

\item \textbf{Change propagation:} a change to the shared abstraction can affect every caller.

\item \textbf{Release coordination:} independently evolving modules may become tied to one shared implementation.

\item \textbf{Conditional complexity:} a shared function may accumulate flags or branches to support callers with different needs.

\item \textbf{Semantic coupling:} two business concepts become connected because they currently use the same implementation.

\item \textbf{Reduced local clarity:} a reader must navigate to a generic component to understand a simple operation.

\end{itemize}

Fowler explains that coupling is not something that can or should be eliminated completely. Modules must communicate to form a functioning system. The practical goal is to control unnecessary or harmful dependencies, particularly where a change in one component forces unrelated components to change \cite{fowler2001coupling}.

\subsection{Premature abstraction}

\textbf{Premature abstraction} occurs when a developer creates a general-purpose interface or framework before the common behavior and its future requirements are understood well enough.

A typical sequence is:

\begin{enumerate}

\item A developer sees two similar code fragments.

\item The developer assumes they represent one stable concept.

\item A generic interface or framework is designed.

\item New parameters, flags, and branches are added as the use cases diverge.

\item Callers become more difficult to understand than the original duplicated code.

\end{enumerate}

This can be worse than keeping two small, explicit implementations. The abstraction may save a few lines but introduce several concepts that every maintainer must understand.

\subsection{Example of overusing DRY}

The file \path{src/dry/premature_abstraction.cpp} creates:

\begin{itemize}

\item a PercentagePolicy interface;

\item separate tax and discount policy classes;

\item a generic PercentageEngine; and

\item a helper that connects the policy to the engine.

\end{itemize}

For this particular example, each policy only multiplies an amount by a fixed percentage. If the requirements are this small and the policies do not need runtime substitution, the abstraction is arguably more complicated than two clearly named calculations.

This does not mean that a policy interface is always wrong. It could be justified if policies are loaded dynamically, supplied by plugins, configured at runtime, or share a larger stable contract. The design is excessive only relative to the current requirements shown in this example.

\subsection{When duplication can be acceptable}

Keeping similar code separate can be safer when:

\begin{itemize}

\item the fragments represent different business policies;

\item the policies have different owners or approval processes;

\item they are expected to change independently;

\item a shared abstraction would need many flags or special cases;

\item the duplicated implementation is small and stable; or

\item the commonality is not yet understood well enough to name and model correctly.

\end{itemize}

This is sometimes described informally as preferring a little duplication over the wrong abstraction. It should not become an excuse for careless copying. Separate implementations should have clear names, explicit ownership, and appropriate tests.

\subsection{Example: independent tax and discount policies}

In \path{src/dry/independent_rules.cpp}, the tax calculator and loyalty-discount calculator both multiply an amount by 10 percent. The shared arithmetic is intentionally simple, but the two calculations represent separate business concepts.

The tax rate may be determined by legal or regulatory requirements. A loyalty discount may be determined by a commercial policy. A change to one does not necessarily imply a change to the other.

Keeping the methods separate makes their meaning explicit and prevents a shared business-policy abstraction from implying that the rates must evolve together. If a stable lower-level arithmetic utility later becomes valuable, it can be introduced without merging the policy ownership.

\textbf{Trade-off:} a few repeated arithmetic characters are acceptable when they preserve the independence of two business rules. The decision should be revisited if the policies later become genuinely governed by one shared rule.

\section{SRP: Does every responsibility require its own class?}

No. SRP does not mean one class must contain only one method. A class can have several methods and still have one clear responsibility if those methods support the same coherent concept and tend to change together.

\subsection{A class can have many methods and one responsibility}

Consider a StudentRecord class with methods to:

\begin{itemize}

\item update the student's name;

\item add a grade;

\item calculate the average grade; and

\item determine academic standing.

\end{itemize}

All of these methods relate to maintaining and interpreting a student's academic record. They can reasonably belong together.

The relevant question is not "How many methods are in this class?" It is "Do these methods serve the same responsibility, or are they affected by unrelated change requests?"

\subsection{Problems caused by too many small classes}

Over-applying SRP may create one class for every getter, setter, or tiny operation. This can produce:

\begin{itemize}

\item \textbf{Excessive indirection:} a simple operation requires navigating through several files and objects.

\item \textbf{More dependencies:} callers must construct or inject many objects for one small task.

\item \textbf{Poor discoverability:} related behavior is scattered across a large number of names.

\item \textbf{Boilerplate:} constructors, interfaces, wrappers, and forwarding methods add code without meaningful behavior.

\item \textbf{False separation:} classes that cannot realistically change or operate independently are treated as if they were independent components.

\end{itemize}

A large number of classes is not automatically a sign of good modularity. A useful class boundary should represent a meaningful responsibility.

\subsection{Interfaces and wrappers are not free}

Interfaces, adapters, wrappers, and intermediate layers are valuable when they provide a real benefit, such as substitutability, dependency inversion, isolation from an external system, or a stable contract.

They are not automatically beneficial when used only to move a one-line getter into another file. Every layer introduces a concept and often a dependency. The benefit of an abstraction should be greater than the mental and maintenance cost it adds.

\subsection{Cohesion and coupling}

Two useful design concepts help evaluate whether a class should be split:

\textbf{Cohesion} describes how strongly the responsibilities and operations inside a module belong together. A highly cohesive component has a clear purpose; its methods contribute to that purpose. A component that combines unrelated tasks has weaker cohesion.

\textbf{Coupling} describes dependencies between components. Some coupling is necessary: components must collaborate to form a system. The concern is unnecessary or fragile coupling, where a change in one component forces changes in others without a strong domain reason.

A useful design often aims for:

\begin{itemize}

\item high cohesion within a component;

\item controlled dependencies between components;

\item explicit interfaces where they provide a real boundary; and

\item no artificial separation that makes normal collaboration harder.

\end{itemize}

Fowler's discussion of coupling emphasizes that dependency is not inherently bad. The pattern and consequences of dependencies matter \cite{fowler2001coupling}.

\subsection{Example: grouping student-record functions}

\path{src/srp/student_cohesive.cpp} groups student-record operations in one class. The class manages the record's name and grades and derives the average and academic standing from that information.

The example deliberately keeps console presentation in main() rather than making the record class responsible for formatting and printing a report. This illustrates a useful boundary: domain data and behavior can stay together, while presentation can remain outside the domain model.

The class has several methods, but they contribute to one coherent responsibility. Splitting it into a separate class for each method would make the design harder to follow without providing a meaningful independent change boundary.

\subsection{Example: excessive SRP decomposition}

\path{src/srp/over_split_student.cpp} defines separate classes for:

\begin{itemize}

\item reading the student's name;

\item setting the student's name;

\item reading the student's age; and

\item setting the student's age.

\end{itemize}

In the shown context, these classes only wrap basic field access. They do not own independent policies or complex behavior. A single Student or StudentRecord type would be clearer.

If age changes later require validation, authorization, or audit logging, those behaviors may justify a meaningful operation or policy. They do not automatically justify one class for every read and write operation.

\textbf{Trade-off:} keep related record operations together until there is a real reason for separate ownership, independent variation, or isolation.

\section{Personal Design Evaluation and Decision-Making}

\subsection{Are DRY and SRP mandatory rules or guiding principles?}

They are \textbf{guiding design principles}, not rules that should be followed blindly.

Their purpose is to help developers reason about maintainability. The principles do not know the complete context of a particular system. A developer must still consider domain meaning, likely change, team ownership, technical constraints, testing, and the cost of introducing abstractions.

A literal interpretation can be harmful:

\begin{itemize}

\item "Never repeat code" can lead to abstractions that unite unrelated policies.

\item "One responsibility per class" can lead to a class for every method.

\item "No coupling" would prevent components from collaborating at all.

\end{itemize}

A better interpretation is to avoid duplicated knowledge and unrelated reasons to change while keeping the design understandable.

\subsection{How to decide whether refactoring is worthwhile}

Refactoring is worthwhile when it improves the internal design enough to justify the change and its future maintenance cost.

Martin Fowler defines refactoring as changing the internal structure of software to make it easier to understand and cheaper to modify without changing its observable behavior \cite{fowler2004refactoring}. His description emphasizes a sequence of small, behavior-preserving transformations rather than a large, risky rewrite.

Before refactoring, consider:

\begin{table}[htbp]
\centering
\small
\begin{tabularx}{\textwidth}{>{\raggedright\arraybackslash}p{0.27\textwidth}X}
\toprule
Factor & Questions to ask \\
\midrule
Meaning & Do the fragments represent the same knowledge or only similar syntax? \\
Change pattern & Are the parts likely to change for the same reason? \\
Ownership & Are they governed by the same business policy or stakeholder? \\
Stability & Is the common behavior understood and stable, or still evolving? \\
Abstraction cost & Will the new API be simpler than the duplicated code? \\
Coupling & Will sharing the code connect previously independent modules? \\
Testability & Can behavior be checked before and after the refactoring? \\
Scale & Is the problem frequent or important enough to justify the new structure? \\
Team comprehension & Will another developer understand where to make the next change? \\
\bottomrule
\end{tabularx}
\end{table}

A refactoring should have a concrete expected benefit, such as removing inconsistent policy copies, isolating an external integration, or making a module easier to test. Refactoring solely to make a diagram look more object-oriented is not a sufficient reason.

\subsection{What to consider before extracting shared code}

Before extracting a common function or module:

\begin{enumerate}

\item Confirm that the fragments have the same semantic purpose.

\item Identify whether they must obey the same policy over time.

\item Compare their edge cases and error behavior.

\item Give the shared concept a meaningful name.

\item Check whether the abstraction needs many flags or special cases.

\item Consider whether callers should depend on the shared component.

\item Keep independent policy decisions outside the shared implementation.

\item Add tests that express the common contract.

\end{enumerate}

A good shared abstraction represents a real commonality. It should not force separate business concepts into one interface merely because they currently use the same expression.

\subsection{What to consider before splitting a class}

Before dividing a class:

\begin{enumerate}

\item Identify the different groups of behavior.

\item Determine what kinds of requests can cause each group to change.

\item Check whether the groups share state and invariants.

\item Assess whether the proposed classes can be understood and tested independently.

\item Examine the dependencies that the split will create.

\item Avoid classes that only forward calls without adding a clear responsibility.

\item Keep the resulting design navigable and understandable.

\end{enumerate}

A class should be split when the new components have meaningful purposes and the separation reduces the impact of independent changes. A class should not be split merely because it contains multiple methods.

\subsection{A situation in which I would deliberately avoid over-applying DRY or SRP}

I would keep two simple calculations separate if they currently use the same formula but represent independent business rules—for example, tax and loyalty discount rates.

The reason is that the policies have different sources of change. A tax regulation can change without any reason for the loyalty program to change. A shared policy object could introduce a false relationship between them. In this situation, separate functions with domain-specific names are easier to understand and safer to evolve.

Likewise, I would not create four classes for reading and writing a student's name and age. Those operations are simple parts of one student-record concept. A single cohesive class provides a clearer place to find and maintain them.

These are deliberate decisions, not an argument against DRY or SRP. They apply the principles according to the meaning and expected evolution of the software.

\section{Code Review: Evaluation of the Five Scenarios}

The assignment asks whether each situation is a violation, a potential problem, or acceptable. The decisions below are based on the context given. Where the scenario does not provide enough implementation detail, the conclusion is conditional rather than absolute.

\subsection{Shared validation}

\textbf{Scenario:} Two operations contain the same customer-validation logic and follow the same business policy.

\textbf{Decision: Share the validation logic.}

This is a strong DRY candidate because the two implementations represent the same knowledge. If the policy changes, both operations should follow the new rule. A single validation function or policy component reduces the risk that one operation becomes inconsistent with the other.

The common component should express the shared rule clearly. If one operation has additional requirements, those should remain explicit rather than being hidden behind a large collection of boolean flags.

\textbf{Classification:} duplicated knowledge worth consolidating.

\subsection{Similar calculations}

\textbf{Scenario:} A tax module and a discount module currently contain the same formula, but their rules come from different business policies.

\textbf{Decision: Keep the policies separate unless a genuine shared rule is established.}

The same arithmetic does not prove the same business meaning. Tax and discount policies can have different legal, commercial, approval, and testing requirements. A future change to tax should not automatically affect a loyalty discount.

The modules may use a low-level mathematical utility if that utility is truly independent of the business policies. However, their policy definitions and domain-specific tests should remain separate.

\textbf{Classification:} acceptable duplication when it preserves independent business rules.

\subsection{A Student class with many methods}

\textbf{Scenario:} A Student class updates student information, calculates GPA, checks academic standing, and displays a summary.

\textbf{Decision: Potential SRP issue; inspect the responsibility behind each method.}

The number of methods is not the deciding factor. Updating a student's record, adding grades, and calculating values derived from those grades may form a cohesive student-record responsibility.

However, academic policy and presentation can have different reasons to change. GPA rules may be controlled by academic regulations, while summary formatting may be requested by a user-interface or reporting stakeholder. If the class directly handles all formatting and output as well as maintaining the record, it may have multiple responsibilities.

A reasonable design is to keep the student data and closely related behavior together, put independent academic policies in a suitable domain service when necessary, and let a presentation or reporting component format the summary.

\textbf{Classification:} potential SRP violation, depending on whether the methods belong to one cohesive responsibility.

\subsection{Too many classes}

\textbf{Scenario:} A developer creates StudentNameGetter, StudentNameSetter, StudentAgeGetter, and StudentAgeSetter.

\textbf{Decision: This is excessive decomposition in the described context.}

Each class wraps one elementary read or write operation. No meaningful independent policy, ownership, or variation is described. The separation increases the number of objects and files that must be understood without providing a clear benefit.

A student-record class can expose appropriate accessors and update operations. If a particular field later gains complex validation or authorization rules, that behavior can be modeled at the point where the policy genuinely belongs.

\textbf{Classification:} over-application of SRP.

\subsection{Order service}

\textbf{Scenario:} An OrderService validates an order, calculates the total, saves the order, and sends a confirmation email.

\textbf{Decision: Separate the independent concerns, while retaining orchestration if useful.}

Validation, pricing, persistence, and notification are likely to change for different reasons. Assigning them to separate components allows each to be tested and modified without embedding all implementation details in one service.

The service may remain as an orchestrator. Its responsibility is to coordinate the order-placement workflow and ensure that the steps occur in the correct order. It should not need to know how a database stores an order or how an email provider sends a message.

In a very small application, a separate class for every operation may be unnecessary. The decision should depend on the complexity, likely change, and need to replace or test each concern independently.

\textbf{Classification:} likely multiple responsibilities; separation is justified when the boundaries represent real change axes.

\subsection{Summary decision table}

\begin{table}[htbp]
\centering
\small
\begin{tabularx}{\textwidth}{>{\raggedright\arraybackslash}p{0.22\textwidth}XXX}
\toprule
\# & Scenario & Assessment & Recommended action \\
\midrule
1 & Shared customer validation & Duplicated knowledge & Extract one shared policy or function \\
2 & Similar tax and discount calculations & Independent business rules & Keep policies separate; share only truly common low-level behavior \\
3 & Student with several methods & Potential SRP problem & Group record behavior; separate independent academic-policy or presentation concerns \\
4 & Separate getter/setter classes & Excessive decomposition & Use a cohesive student-record type \\
5 & Multi-purpose order service & Likely multiple reasons to change & Separate validator, calculator, repository, and sender; keep orchestration if it has a clear purpose \\
\bottomrule
\end{tabularx}
\end{table}

\section{A Practical Decision Framework}

The following framework turns the principles into a repeatable code-review process.

\subsection{When considering DRY}

Ask:

\begin{enumerate}

\item Is the repeated fragment the representation of one shared fact, rule, or policy?

\item Would a change to that knowledge need to update all copies together?

\item Does a shared abstraction have a clear name and stable contract?

\item Would the abstraction create a dependency between concepts that should vary independently?

\item Is the cost of maintaining the copies greater than the cost of maintaining the abstraction?

\end{enumerate}

If the first two answers are yes and the abstraction is simple, sharing is likely useful. If the fragments only look alike and their policies change independently, keeping them separate may be safer.

\subsection{When considering SRP}

Ask:

\begin{enumerate}

\item What responsibility does the component represent?

\item Which stakeholder, policy, or technical concern can request a change?

\item Do groups of methods change for different reasons?

\item Do the methods share important state or invariants?

\item Would splitting the component improve cohesion and reduce harmful coupling?

\item Would the split add more indirection than value?

\end{enumerate}

A useful split creates components that can be understood and changed for distinct reasons. A useless split creates small wrappers with no independent meaning.

\subsection{Safe refactoring sequence}

A practical refactoring sequence is:

\begin{enumerate}

\item \textbf{Observe current behavior.} Write tests or record representative input/output cases.

\item \textbf{Identify the design issue.} State which knowledge is duplicated or which responsibilities are mixed.

\item \textbf{Make one small structural change.} Avoid combining several unrelated refactorings.

\item \textbf{Run tests.} Confirm that externally visible behavior remains the same.

\item \textbf{Review dependencies.} Check whether the new abstraction creates coupling or unnecessary indirection.

\item \textbf{Evaluate readability.} Ask whether another developer can more easily locate the next change.

\item \textbf{Stop when the design is clear enough.} Do not continue splitting or generalizing without a concrete benefit.

\end{enumerate}

This sequence follows the behavior-preserving approach to refactoring described by Fowler \cite{fowler2004refactoring}.

\section{Conclusion}

DRY and SRP address different but related design problems. DRY discourages maintaining the same knowledge in several places. SRP encourages grouping behavior according to coherent responsibilities and reasons to change. Both can improve maintainability when their underlying concepts are understood.

However, neither principle should be interpreted literally as a demand to eliminate every repeated expression or to create one class for every method. Superficially similar code may implement independent policies. A class may contain many methods that all support one responsibility. Shared abstractions can introduce dependencies, and excessive class decomposition can make code harder to navigate.

The examples and code-review scenarios demonstrate that good design requires judgment. Developers should consider semantic meaning, policy ownership, expected change, cohesion, coupling, testability, and the cost of abstraction. A little duplication can be safer than an incorrect abstraction, and a cohesive class can be clearer than many tiny classes.

The most useful objective is not minimum duplication or maximum decomposition. It is a codebase that is understandable, testable, and reasonably easy to change without introducing unnecessary complexity.

\printbibliography[heading=bibintoc,title={References}]

\appendix
\section{DRY Example: Duplicated Customer Validation}
\lstinputlisting[caption={Duplicated validation before refactoring},label={lst:validation-before}]{../src/dry/validation_before.cpp}

\section{DRY Example: Shared Customer Validation}
\lstinputlisting[caption={Shared validation after refactoring},label={lst:validation-after}]{../src/dry/validation_after.cpp}

\section{DRY Example: Independent Business Rules}
\lstinputlisting[caption={Similar calculations representing independent policies},label={lst:independent-rules}]{../src/dry/independent_rules.cpp}

\section{DRY Example: Premature Abstraction}
\lstinputlisting[caption={An over-generalized percentage abstraction},label={lst:premature-abstraction}]{../src/dry/premature_abstraction.cpp}

\section{SRP Example: Multi-Responsibility Order Service}
\lstinputlisting[caption={Order service before separation of responsibilities},label={lst:order-before}]{../src/srp/order_service_before.cpp}

\section{SRP Example: Separated Order Responsibilities}
\lstinputlisting[caption={Order service after separation of responsibilities},label={lst:order-after}]{../src/srp/order_service_after.cpp}

\section{SRP Example: Cohesive Student Record}
\lstinputlisting[caption={A cohesive student-record class with several related methods},label={lst:student-cohesive}]{../src/srp/student_cohesive.cpp}

\section{SRP Example: Excessive Student Class Decomposition}
\lstinputlisting[caption={Separate getter and setter classes with little independent responsibility},label={lst:student-over-split}]{../src/srp/over_split_student.cpp}

\end{document}
